\documentclass[11pt]{article}
\usepackage{mathblatt}
\makeatletter
\setlength{\mbpfbe}{0mm}% keine Punkte (Bauregel 6.4)
\makeatother
% Rechenraum als Karo (Bauregel 6.9): n Rechenschritte = n mal 10 mm
\newcommand{\karo}[1]{\par\smallskip\noindent\begin{tikzpicture}
  \clip (0,0) rectangle ({\linewidth-1mm},{#1*10mm});
  \draw[black!16,line width=0.3pt,step=5mm] (0,0) grid ({\linewidth},{#1*10mm});
  \end{tikzpicture}\par}
\newcommand{\linie}[1]{\,{\color{black!45}\rule[-1pt]{#1}{0.4pt}}\,}
\newcommand{\pkt}{\textperiodcentered{}}
\begin{document}

\pfheftstil
\pflaufkopf[P10 \pkt{} mehr]{Kenngrößen}{PS30-3}\pfseitevon
\pfheftkopf{Kenngrößen}{mehr}
\pfstufenzeile{Mittelwert \pkt{} Minimum, Maximum, Spannweite \pkt{} Median}
% 2021-OS-B1f-m1 (Schwere I, rang 11)
\begin{pfaufg}{nach~P10\\’21~\pkt{}~1f}{1.}{}
Jonas springt im Weitsprung 4,24 m, 3,96 m, 4,10 m und 4,34 m. Gib seine durchschnittliche Weite an.
\karo{2}
\fusshilfe{\mbox{1:~4,16 m}}
\end{pfaufg}
% 2020-OS-K2c-m1 (Schwere I, rang 12)
\begin{pfaufg}{nach~P10\\’20~\pkt{}~2c}{2.}{}
Ein Schwimmbad zählt in der ersten Woche der Sommerferien jeden Tag seine Besucher. Am Donnerstag kamen 155 Besucher. Im Balkendiagramm fehlt der Balken für Donnerstag noch. Berechne, wie viele Besucher in dieser Woche durchschnittlich an einem Tag im Schwimmbad waren.
\par\smallskip\noindent{\small\itshape Skizze: Balkendiagramm (Achse beginnt bei 0), Anzahl der Besucher: Mo: 150, Di: 120, Mi: 180, Do: leer, Fr: 170, Sa: 210, So: 240}\par
\karo{2}
\fusshilfe{\mbox{2:~175 Besucher}}
\end{pfaufg}
% 2019-OS-B1i-m1 (Schwere I, rang 13)
\begin{pfaufg}{nach~P10\\’19~\pkt{}~1i}{3.}{}
Gib die Spannweite dieser Daten an: 8,4 m; 8,9 m; 8,6 m; 8,2 m; 8,7 m; 8,4 m
\karo{2}
\fusshilfe{\mbox{3:~0,7 m}}
\end{pfaufg}
% 2017-OS-B1h-m1 (Schwere I, rang 14)
\begin{pfaufg}{nach~P10\\’17~\pkt{}~1h}{4.}{}
Gib den Durchschnitt (das arithmetische Mittel) dieser Längen an: 2,4 m; 2,2 m; 2,0 m; 2,2 m
\karo{2}
\fusshilfe{\mbox{4:~2,2 m}}
\end{pfaufg}
% 2016-OS-B1a-m1 (Schwere I, rang 15)
\begin{pfaufg}{nach~P10\\’16~\pkt{}~1a}{5.}{}
Gib das arithmetische Mittel (den Durchschnitt) der Werte 12; 45; 60 an.
\karo{2}
\fusshilfe{\mbox{5:~39}}
\end{pfaufg}
% 2015-OS-B1f-m1 (Schwere I, rang 16)
\begin{pfaufg}{nach~P10\\’15~\pkt{}~1f}{6.}{}
Gib den Median (Zentralwert) dieser Messwerte an: 6 °C; 2 °C; 9 °C; 4 °C; 9 °C; 7 °C; 4 °C
\karo{2}
\fusshilfe{\mbox{6:~6 °C}}
\end{pfaufg}
% 2015-OS-K7a-m1 (Schwere I, rang 17)
\begin{pfaufg}{nach~P10\\’15~\pkt{}~7a}{7.}{}
Eine Fußballmannschaft hatte in einer Saison 18 Heimspiele. Zu diesen Spielen kamen zusammen 812 467 Zuschauer. Berechne, wie viele Zuschauer im Durchschnitt bei einem Heimspiel waren. Runde sinnvoll.
\karo{2}
\fusshilfe{\mbox{7:~45137 Zuschauer}}
\end{pfaufg}
% 2014-OS-K4a-m1 (Schwere I, rang 18)
\begin{pfaufg}{nach~P10\\’14~\pkt{}~4a}{8.}{}
Eine Tankstelle hat im Sommer 2013 zehn Wochen lang einmal pro Woche die Preise für die Kraftstoffe E 10 und Diesel notiert (in Cent pro Liter). Lies aus der Tabelle den niedrigsten und den höchsten Preis für E 10 ab (Minimum und Maximum).
\par\smallskip\noindent{\small\itshape Skizze: Tabelle: Datum | 01.07. | 08.07. | 15.07. | 22.07. | 29.07. | 05.08. | 12.08. | 19.08. | 26.08. | 02.09.; E 10 | 151,8 | 153,6 | 156,3 | 157,0 | 158,5 | 157,9 | 160,4 | 165,1 | 166,3 | 163,8; Diesel | 137,2 | 139,9 | 141,0 | 143,3 | 142,5 | 144,8 | 147,6 | 151,4 | 150,2 | 149,0}\par
\karo{3}
\fusshilfe{\mbox{8:~Minimum am 01.07.: 151,8 ct; Maximum am 26.08.: 166,3 ct}}
\end{pfaufg}
% 2014-OS-K4b-m1 (Schwere I, rang 19)
\begin{pfaufg}{nach~P10\\’14~\pkt{}~4b}{9.}{}
Eine Tankstelle hat im Sommer 2013 zehn Wochen lang einmal pro Woche die Preise für die Kraftstoffe E 10 und Diesel notiert (in Cent pro Liter). Berechne die Spannweite der Dieselpreise in diesen zehn Wochen.
\par\smallskip\noindent{\small\itshape Skizze: Tabelle: Datum | 01.07. | 08.07. | 15.07. | 22.07. | 29.07. | 05.08. | 12.08. | 19.08. | 26.08. | 02.09.; E 10 | 151,8 | 153,6 | 156,3 | 157,0 | 158,5 | 157,9 | 160,4 | 165,1 | 166,3 | 163,8; Diesel | 137,2 | 139,9 | 141,0 | 143,3 | 142,5 | 144,8 | 147,6 | 151,4 | 150,2 | 149,0}\par
\karo{2}
\fusshilfe{\mbox{9:~14,2 ct}}
\end{pfaufg}
\par\vfill\noindent{\scriptsize\color{mbgrau}Originale: P10 2021 \pkt{} 1f, P10 2020 \pkt{} 2c, P10 2019 \pkt{} 1i, P10 2017 \pkt{} 1h, P10 2016 \pkt{} 1a, P10 2015 \pkt{} 1f, P10 2015 \pkt{} 7a, P10 2014 \pkt{} 4a, P10 2014 \pkt{} 4b\par}
\end{document}